\subsection{Elementary optimisation models}

An optimisation problem has two parts: building the correct function and analyzing it on the allowed domain.

\begin{examplebox}[Maximum area with fixed perimeter]
A rectangle has perimeter \(20\). If its width is \(x\), then its length is \(10-x\), so
\[
A(x)=x(10-x)=10x-x^2,
\qquad 0<x<10.
\]
We compute
\[
A'(x)=10-2x.
\]
The critical point is \(x=5\). Since
\[
A''(x)=-2<0,
\]
this point gives a maximum. The optimal rectangle is a square of side \(5\), with area \(25\).
\end{examplebox}

\begin{examplebox}[Minimum material in a simple open box model]
An open box with square base has volume \(32\). If the base side is \(x>0\), then the height is \(h=32/x^2\). Its surface area is
\[
S(x)=x^2+4xh=x^2+\frac{128}{x}.
\]
Thus
\[
S'(x)=2x-\frac{128}{x^2}.
\]
The condition \(S'(x)=0\) gives \(2x^3=128\), hence \(x=4\). The corresponding height is \(h=2\).
\end{examplebox}

\begin{interpretationbox}[Why the model matters more than the derivative]
The differentiation is elementary. The substantial reasoning lies in choosing a variable, expressing all quantities through it, identifying the admissible domain, and interpreting the critical point in the original problem.
\end{interpretationbox}
